%------------------------------------------------------------------------------%
\documentclass[fleqn,utf8,aspectratio=169,12pt]{beamer}

\usepackage{integracao_e_series}

\title{Teorema do Valor Médio para Integrais}

%------------------------------------------------------------------------------%
\begin{document}

\addtitle

%------------------------------------------------------------------------------%
\section{Valor Médio}

%------------------------------------------------------------------------------%
\begin{frame}{Valor Médio}
  O \emph{valor médio} de $f$ em $[a, b]$ é
  \pause
  o valor que produz a mesma área que a função
  \[
    \pause
    \fmedio\;(b-a) = \int_a^b f(x) dx
  \]

  \pause
  Podemos calcular o valor médio a partir da integral definida
  \[
    \pause
    \fmedio = \frac{1}{b-a} \int_a^b f(x) dx
  \]
\end{frame}

%------------------------------------------------------------------------------%
\section{Teorema do Valor Médio para Integrais}

%------------------------------------------------------------------------------%
\begin{frame}{Valor Médio}
  \centering
  \bigskip
  \input{fig_valor_medio_media}
\end{frame}

%------------------------------------------------------------------------------%
\begin{frame}{Valor Médio}
  O Teorema do Valor Médio garante que\\[1mm]
  \pause
  \medskip
  \(\fmedio\) sempre é o valor de $f$ para algum ponto em $[a, b]$
\end{frame}

%------------------------------------------------------------------------------%
\begin{frame}{Teorema do Valor Médio para Integrais}
  Seja $f$ uma função {\color{blue} contínua} em $[a, b]$, \pause
  então existe ${\color{blue} c} \in (a, b)$ \pause
  tal que
  \[
    f({\color{blue}c})(b-a)
    \;=\; \int_a^b f(x)\,dx
  \]
  \pause
  ou
  \[
    f({\color{blue}c})
    \;=\; \frac{1}{b-a}\int_a^b f(x)\,dx
  \]
\end{frame}

%------------------------------------------------------------------------------%
\section{Demonstração}

%------------------------------------------------------------------------------%
\begin{proofframe}{Demonstração}
  Queremos mostrar que existe $c \in (a, b)$ tal que
  \[
    f(c) = \frac{1}{b-a}\int_a^b f(x)\,dx
  \]
\end{proofframe}

%------------------------------------------------------------------------------%
\begin{proofframe}{Demonstração}
  Teorema de Weierstrass \\
  \pause
  \qquad $f$ é \emph{contínua} em $[a, b]$
  \pause
  então existe um \emph{mínimo} $f(u)$ e um \emph{máximo} $f(v)$

  \pause
  {\centering
  \bigskip
  \onslide<+->{\input{fig_valor_medio_media}}}
\end{proofframe}

%------------------------------------------------------------------------------%
\begin{proofframe}{Demonstração}
  \onslide<+->{Como a integral \emph{preserva a desigualdade}}
  \begin{align*}
    \onslide<+->{\int_a^b f(u)\,dx \;\leq\; & \int_a^b f(x)\,dx \;\leq\; \int_a^b f(v) \,dx \\[3mm]}
    \onslide<+->{f(u)(b-a) \;\leq\;         & \int_a^b f(x)\, dx \;\leq \;f(v)(b-a)         \\[3mm]}
    \onslide<+->{f(u) \;\leq\;               \frac{1}{b-a} & \int_a^b f(x)\,dx \;\leq\; f(v)}
  \end{align*}
\end{proofframe}

%------------------------------------------------------------------------------%
\begin{proofframe}{Demonstração}
  Valor médio
  \[
    \fmedio = \frac{1}{b-a} \int_a^b f(x)\,dx
  \]
  \pause
  \[
    f(u) \;\leq\; f_{\text{médio}} \;\leq\; f(v)
  \]

  \pause
  \bigskip
  Teorema do Valor Intermediário

  \pause
  \qquad $f$ precisa assumir o valor \(\fmedio\) em algum $c\in[a, b]$
\end{proofframe}

%------------------------------------------------------------------------------%
\section{Lista Mínima}

%------------------------------------------------------------------------------%
\begin{frame}{Lista Mínima}
  Estudar a Seção 2.5 da Apostila

  \vfill
  Atenção: A prova é baseada no livro, não nas apresentações
\end{frame}

%------------------------------------------------------------------------------%
\end{document}
%------------------------------------------------------------------------------%
